\subsection{Mixed variations and a disappearing exit}
\label{rw:app:relative-variations}

The proof of Proposition~\ref{rw:relative-core} has two parts.  Normal
variational equations give decay of mixed jets at fixed time; the
fixed transverse exit then controls the hitting-time jets.  Normal
decay dominates their growth, so differentiated endpoint terms vanish
at a lost exit and the relative integral extends to infinite time.
We work on one common real domain and choose the derivative order
before reducing it.  Retaining $N+2$ joint derivatives of the
triangular field, initial sections and coefficient changes gives
relative jets through total order $N$.  Any two derivatives spent
later on a logarithmic rate must already be included in $N$ here.

Let $w(t;a,\mathbf b)$ be the scalar solution starting at $a$ and
let $Z(t;z,a,\mathbf b)$ be its normal solution. A selected derivative
means a product of ordinary $a,z,\mathbf b$ derivatives and, in the
lower case, $E_\epsilon$. Euler derivatives obey the same product
rule as ordinary derivatives; when commuting them through the vector
field, derivatives of $q=\epsilon$ retain its factor $\epsilon$.
In the upper case $q=1$.  Differentiation in frozen parameters gives
inhomogeneous variational equations and hence polynomial factors in
time.  We estimate these equations directly, using the scalar-flow
bounds.

For a normal jet of positive total order, differentiating
$\dot Z=-g(w,Z)$ gives
\begin{equation}\label{rw:app:normal-induction}
 \frac{\dd}{\dd t}D^\alpha Z
 =-g_Z(w,Z)D^\alpha Z+\mathcal R_\alpha(t).
\end{equation}
Terms in $\mathcal R_\alpha$ containing only scalar and coefficient
variations still contain $Z$: all derivatives of $g(w,0)$ in those
variables vanish. Every other term contains a lower normal variation.
The scalar variations satisfy
\eqref{rw:relative:slow-growth}. At order one, variation of constants
in \eqref{rw:app:normal-induction} gives a polynomial time factor
times $e^{-(b_*-c_1\varpi)t}$. At order $k$, products of previously
bounded variations have the same form with a larger polynomial and
a larger constant multiplying $\varpi$. The convolution estimate
\[
 \int_0^t e^{-b_*(t-s)}(1+s)^m
                 e^{-(b_*-c\varpi)s}\,\dd s
 \le C(1+t)^{m+1}e^{-(b_*-c\varpi)t}
\]
closes the induction. Include the finitely many constants needed
through order $N+1$ in $c_N$ before choosing $\varpi$. We obtain
\begin{equation}\label{rw:app:normal-jets}
 |D^\alpha Z(t)|\le
 C_N(1+t)^{M_N}e^{-(b_*-c_N\varpi)t},\qquad |\alpha|\le N+1.
\end{equation}
For $\alpha=0$, the stronger estimate
$|Z(t)|\le|z|e^{-b_*t}$ holds. At normal input zero, normal-input
variations in \eqref{rw:app:normal-jets} need not vanish; they still
decay.  Their decay makes the relative integral differentiable at zero
normal input as well.

The identity
\[
 \beta(w,Z)-\beta(w,0)
 =Z\int_0^1\beta_Z(w,\theta Z)\,\dd\theta
\]
and \eqref{rw:app:normal-jets} now bound every derivative of its
left-hand side by \eqref{rw:relative:majorant}. Replacing $\beta$
by $g_Z$ proves the corresponding estimate for the logarithm of the
normal derivative.  The estimate is uniform through $z=0$.

Write the fixed-exit equation as
$H(T,\theta)=\phi_T(a(\theta);\mathbf b(\theta))-v_R(\theta)=0$.
Its time derivative has magnitude at least $cq$.
At order $k$, the chain rule separates $H_tD^\alpha T$ from
products
$(\partial_t^mD^\gamma H)(T)\prod_{\ell=1}^mD^{\beta_\ell}T$,
where $|\gamma|+\sum|\beta_\ell|=k$ and every remaining time jet
has order below $k$. Mixed time derivatives of the flow follow
from its equation with the same finite-order slow-growth bounds.
Dividing by the fixed-exit velocity proves
\eqref{rw:relative:time-jets} by induction. Regular parameter
dependence of either section is included in $H$. The identity
$\partial_a\phi_T=f(v_R)/f(a)$ already places any reciprocal
entry speed in the scalar variational bound; no entry-speed margin
is required.

For a quadratic passage from $-v_R$ to $v_R$, the more explicit
calculation is useful. With $f(v)=\mu+v^2$ and $\mu>0$,
\[
 T=q^{-1}\int_{-v_R}^{v_R}\frac{\dd v}{\mu+v^2},\qquad
 |\partial_\mu^jT|\le C_jq^{-1}\mu^{-j-1/2}
 \le C_jq^{2j}(1+T)^{2j+1}.
\]
Thus a vanishing interior gap costs only a fixed power of the transit
time in this case. At a split root, the preceding implicit-flow
argument supplies the appropriate slow exponential bound instead.

At total order $k$, differentiating an integral with endpoint $T$
gives its fixed-time differentiated integral and a finite sum of terms
\begin{equation}\label{rw:app:endpoint-term}
 (\partial_t^{m-1}D^\gamma\mathcal I)(T)
       \prod_{\ell=1}^mD^{\beta_\ell}T,
 \qquad m\ge1,\quad |\beta_\ell|\ge1,\quad
 |\gamma|+\sum_\ell|\beta_\ell|=k.
\end{equation}
The integrand in each product uses total order at most $k-1$.
Its normal factor, the time-jet bounds, and the combined slow exponent
reserved before choosing $\varpi$ bound every term by
$C_Nq^{-M_N}(1+T)^{M_N}e^{-(b_*-2c_N\varpi)T}$.
When $q\to0$, separated sections give $q^{-1}\le T/c$; the
remaining polynomial is absorbed by the positive normal margin.
Thus all endpoint terms vanish uniformly at a disappearing exit.

For the absolute factor $L=e^{-I_0}$, differentiate
$I_0=\int_0^T\beta(w(t),0)\,\dd t$ only at finite $T$.
Its derivatives have polynomial and slow-exponential bounds from the
same flow and time equations. The chain rule for $e^{-I_0}$ then
multiplies finite products of these derivatives by
$L\le e^{-b_*T}$. This proves absolute exponential bounds for all
selected derivatives of $L$.  The action $I_0$ itself may diverge.

At the separated-section slow-root boundaries of
Proposition~\ref{rw:relative-core}, the transit time tends to infinity.
On the complete auxiliary flow define
$j_\infty=-\int_0^\infty\mathcal I(t)\,\dd t$.
The integrable majorant permits differentiation through the integral.
The difference $j_T-j_\infty$ is a tail, and its selected derivatives,
including \eqref{rw:app:endpoint-term}, tend to zero. Define $j=j_T$
on the passing side and $j=j_\infty$ on the nonpassing side. These
derivatives agree at the lost-exit boundary; induction on derivative
order, or the one-variable fundamental theorem along coordinate
segments, proves the resulting $C^N$ extension in the selected ordinary
variables, together with the selected Euler scale bounds. Define the absolute
factor to be zero on the nonpassing side. Its jets tend to zero by
the absolute estimate, giving a separate extension in the same derivative
class.  Both extensions are defined on the prescribed auxiliary domain.
Their restriction to real passing orbits gives the transition maps
used in the comparison.

\subsection{The central kernel on a real coefficient domain}
\label{rw:app:relative-core}

The central connection requires a different estimate: the reference
is controlled with a Gaussian weight rather than by ordinary
high-order state derivatives.  We prove the mixed bounds asserted in
Proposition~\ref{rw:transitions:central-core}.  All integrals may be
rescaled by $X=Ry$ to the fixed interval $0\le y\le1$. Write
$F_\epsilon=F_++\epsilon\bar m$ and split the exact graph field as
\[
 \mathfrak f(X,z)=a_0(X)z+\frac{\nu}{F_\epsilon(X)}+\mathcal N(X,z),
 \qquad a_0=\frac{X}{\epsilon F_\epsilon}-\frac{\bar c}{F_\epsilon}.
\]
Direct differentiation of \eqref{rw:relative:core-field} on its
signed tube gives
\begin{equation}\label{rw:app:core-nonlinearity}
 |\mathcal N_z(X,z)|\le C[(X/\epsilon+1)|z|+|\nu|],
 \qquad \mathcal N(X,0)=0.
\end{equation}
The backward kernel
$K(X,S)=\exp[-\int_X^S a_0(u)\,\dd u]$ satisfies
\begin{equation}\label{rw:app:gaussian-kernel}
 |K(X,S)|\le Ce^{-c(S^2-X^2)/\epsilon},\qquad
 \int_X^R|K(X,S)|\,\dd S
          \le\frac{C\epsilon}{X+\sqrt\epsilon}.
\end{equation}
For the last inequality, split at $X=\sqrt\epsilon$ and use
$S^2-X^2\ge(S-X)^2$ in the first case and
$S^2-X^2\ge X(S-X)$ in the second. The terminal term
$K(X,R)\epsilon\nu h_b$ has the same weight, since $R\ge R_0$.

In the norm
$\|z\|_*:=\sup_X(X+\sqrt\epsilon)|z(X)|/\epsilon$,
the integral of the forcing in the backward equation has size
$C|\nu|$. On $\|z\|_*\le C_0|\nu|$,
$(X/\epsilon)|z|\le C_0|\nu|$; hence
\eqref{rw:app:core-nonlinearity} and
\eqref{rw:app:gaussian-kernel} make its nonlinear part a contraction
after $|\nu|$ and the tube have been reduced. This proves existence,
uniqueness and the first bound in
\eqref{rw:relative:reference-weight} for the continuation of the
given terminal graph. Its uniqueness identifies it with that graph
on the overlap with the lower phase construction.

For each selected coefficient or Euler derivative, differentiating
the kernel introduces at most a fixed polynomial in
$(S^2-X^2)/\epsilon$; the differentiated $R$ endpoint is treated
in the fixed $y$ coordinate. Such polynomials are absorbed by a
smaller positive Gaussian exponent. The terminal data retain their
$\epsilon$ weight, as do their Euler derivatives. Differentiate the
contracting equation. Its highest unknown derivative occurs through
the same integral operator of norm less than $1/2$, while its other
terms involve lower reference jets and the differentiated kernels.
Induction yields the second bound in
\eqref{rw:relative:reference-weight}. The identity $z_b=0$ at
$\nu=0$ and one additional $\nu$ derivative restore its $\nu$
factor whenever no $\nu$ derivative was prescribed. Accordingly,
retain one extra coefficient order when that divided bound is needed.
All of these arguments use real differentiation and the weighted
integral norm.

The linear coefficient about this reference is
$a=\mathfrak f_z(X,z_b)$. The exact identity
\[
 D-zD_z=F_\epsilon-\epsilon\bar m z^2
\]
shows that
$a-X/(\epsilon F_+)$ and its selected coefficient and Euler
jets are bounded pointwise. The term $F_\epsilon-F_+=\epsilon\bar m$
removes the apparent inverse scale in that difference; every remaining
dangerous term contains $(X/\epsilon)z_b$, which the weighted bound
controls. Taylor division of the nonlinear difference gives
\[
 B_c(X,v)=\int_0^1(1-\theta)
           \mathfrak f_{zz}(X,z_b+\theta v)\,\dd\theta.
\]
Its ordinary $v$ derivatives and selected coefficient derivatives
have bounds $C_N(X/\epsilon+1)$ on the fixed signed input tube.
Retaining $N+2$ derivatives of the graph field supplies these $N$
orders; the reference construction may retain $N+3$ when its divided
$\nu$ bound is also used.

The exact kernel $E(X)=\exp[-\int_X^R a]$ has the same Gaussian
bound. Its selected derivatives have that bound times a fixed
polynomial in $(R^2-X^2)/\epsilon$. Consequently,
\begin{equation}\label{rw:app:normalized-kernel-bound}
 \int_0^R (X/\epsilon+1)|D^\alpha E(X)|\,\dd X\le C_N.
\end{equation}
For the principal term, the substitution
$u=(R^2-X^2)/\epsilon$ evaluates the integral with weight
$X/\epsilon$ directly; the remaining term follows from the same
Gaussian estimate and $R\ge R_0$.

Equation \eqref{rw:relative:normalized-volterra} is therefore a
contraction for $Q$ in a fixed small ball and for $W$ in a fixed
interval independent of the scale. Differentiate it successively in
$W$, the coefficients and $E_\epsilon$. The highest $Q$ jet always
appears with its original contracting linearization. All other terms
are integrals of products of lower jets, derivatives of $B_c$, and
at least one kernel factor $E$. Equation
\eqref{rw:app:normalized-kernel-bound} bounds them. Solving at each
order by the inverse of $I-\mathcal T_Q$, whose norm is at most two,
proves the required mixed bounds. This argument works at $W=0$ as
well as at nonzero input. Finally $a_r$ has bounded selected jets,
so its integral supplies the bounds for $\log a_c$ and its positive
upper and lower margins.

These estimates hold at any finite order chosen in advance, with
constants that may depend on that order.  They control ordinary
section and coefficient derivatives and Euler scale derivatives of
the integrated relative functions.  At the Gaussian center the
reference has the stated weighted bound; uniform high-order
$X$ derivatives and ordinary $\epsilon$ derivatives of the final
relative maps remain outside these estimates.
